\begin{center}
\begin{tikzpicture}[x=.88cm,y=.88cm,font=\small]
  \fill[paper] (0,0) rectangle (12.3,6.7);

  \draw[source!70,line width=.8pt] (.55,2.18) .. controls (3.4,1.78) and
    (7.05,2.57) .. (11.82,1.84);
  \node[text=source,below,font=\scriptsize] at (8.3,1.97) {长江方向（示意）};
  \draw[source!55,dashed] (.8,3.03) .. controls (4.0,3.35) and
    (7.85,2.82) .. (11.55,3.38);
  \node[text=source,font=\scriptsize] at (10.73,3.63) {淮河方向（示意）};

  \fill[timeline!24] (.85,1.0) ellipse (1.2 and .68);
  \node[align=center] at (.85,1.0) {巴蜀\\成汉（至347）};

  \fill[dynasty!30] (3.65,1.52) ellipse (1.32 and .83);
  \node[align=center] at (3.65,1.52) {荆、江中游\\东晋军镇};

  \fill[dynasty!44] (8.1,1.03) ellipse (1.58 and .78);
  \node[align=center,text=white] at (8.1,1.03) {建康与江东\\东晋朝廷};

  \fill[source!23] (1.72,4.22) ellipse (1.34 and .9);
  \node[align=center] at (1.72,4.22) {河西\\前凉等政权};

  \fill[timeline!27] (4.58,4.67) ellipse (1.45 and 1.0);
  \node[align=center] at (4.58,4.67) {关中、河东\\前赵—前秦\\及其后继者};

  \fill[source!31] (7.75,4.78) ellipse (1.48 and .98);
  \node[align=center] at (7.75,4.78) {河北、中原\\后赵—冉魏—\\前燕等};

  \fill[timeline!19] (10.52,4.55) ellipse (1.05 and .78);
  \node[align=center,font=\scriptsize] at (10.52,4.55) {辽西、辽东\\慕容诸燕};

  \draw[->,very thick,color=dynasty]
    (7.85,1.8) .. controls (7.55,2.44) and (7.58,2.9) .. (7.68,3.58);
  \node[text=dynasty,right,font=\scriptsize,align=left] at (7.91,2.7)
    {淮北防务与\\北伐方向};

  \draw[->,thick,dashed,color=timeline]
    (4.1,2.18) .. controls (4.42,2.82) and (4.67,3.2) .. (4.72,3.65);
  \node[text=timeline,left,font=\scriptsize,align=right] at (4.0,3.17)
    {荆州—关中\\北伐方向};

  \draw[->,thick,dashed,color=source]
    (6.05,4.63) .. controls (6.38,4.73) and (6.58,4.76) .. (6.8,4.77);
  \node[text=source,above,font=\scriptsize] at (6.45,4.93)
    {北方政权反复更替};

  \node[draw=source!75,fill=paper,rounded corners=2pt,align=center,
    inner sep=3pt,font=\scriptsize] at (10.05,6.0)
    {317--383年间，建康朝廷、地方军镇与北方多政权\\
    在河流、山地与旧城网络之间相互牵动};
\end{tikzpicture}

\smallskip
\small\textit{图\quad 东晋南方朝廷与北方多政权的相对空间，约317--383年：据《晋书·元帝纪》《成帝纪》《孝武帝纪》、〈载记〉及《资治通鉴》卷八十九至一百零五，提示建康、荆州军镇、巴蜀、关中、河北和辽西之间的关系。政权名称按这一长时段先后列示，不表示它们在同一时点同时存在；各色块、河流与箭线均为非等比例的约略示意，不主张精确疆界、治所位置、交通线路或北伐推进线。}
\end{center}

\phantomsection\label{fig:eastern-jin-polities}
